\paragraph{Why the model above cannot reach the apartment wedge.}
Calibration cannot resolve the difficulty. In the model of Section~\ref{ss:model_setup}, a buyer who sees no rating prices the listing at her belief about its quality, and a belief about a binary rating lies between the worst rating and the best. A silent listing therefore sells at most at the price of a verified efficient listing: the ratio of the non-disclosure premium to the energy-efficiency premium is the pool belief, which is bounded by $\bar\theta\le1$ (Lemma~\ref{lem:model_bounds}(i)). The same bound applies to the differences-in-differences whenever the pool belief is stable across the reform.\footnote{The ordering survives exact pricing under the monotonicity condition discussed in Appendix~\ref{app:att_exact}. The levels make the same point: silent apartments sold above verified A-to-D listings by \attLevelGapPreApt{} log points before the reform and by \attLevelGapPostApt{} after, whereas the model prices silence below verified efficiency at any date and for any belief in $[0,1]$. Under Assumption~\ref{assum:no_pre_pricing} the pre-reform gap is non-energy, and the energy component of the post-reform gap is the DiD, \attExcessApt{} log points, still positive.} The apartment wedge of \attLamApt{} log points decomposes accordingly: \attExcessApt{} points are the amount by which $\hat\Delta$ exceeds $\hat\beta_1$, a region the model cannot enter, and only the remaining \attReachApt{} points, $(1-\hat e^*)\hat\beta_1$, are within reach of any change in how buyers form beliefs about silent listings.
